\documentclass[11pt]{article}

% ACL style - two-column, review mode
\usepackage[review]{acl}
\usepackage{times}
% times.sty sets \ttdefault to pcr (Adobe Courier). Some minimal TeX installs
% lack the pcr metric fonts in the active search path; CM typewriter is always
% available under T1 and suffices for \texttt in this paper.
\renewcommand{\ttdefault}{cmtt}
\usepackage{latexsym}
\usepackage[T1]{fontenc}
\usepackage[utf8]{inputenc}
\usepackage{microtype}
\usepackage{graphicx}
\usepackage{booktabs}
\usepackage{amsmath}
\usepackage{amssymb}
\usepackage{multirow}
\usepackage{xcolor}
\usepackage{array}
\usepackage{colortbl}
\usepackage{algorithm}
\usepackage{algpseudocode}

\definecolor{tracebg}{RGB}{236,243,255}
\definecolor{highconf}{RGB}{212,240,212}
\definecolor{lowconf}{RGB}{255,224,224}

\title{Why Ask When You Can Observe?\\
A Vision-Language-Action Model for Epistemic\\
Action Selection in Multi-Agent Crop Disease Diagnosis}

\author{Anonymous Submission \\ EMNLP 2026 Main Conference}

\begin{document}
\maketitle

%-----------------------------------------------------------------------
\begin{abstract}
We show that routing behavior in multi-agent VLM systems---path length,
backtrack decisions, contradiction events---predicts correctness far better
than self-declared confidence. We present \textbf{OBSERVE}, the first
Vision-Language-Action model trained on routing traces from
\textbf{PlantSwarm}, a five-agent diagnostic swarm for plant disease
diagnosis. OBSERVE learns to: (1) predict next-agent routing decisions from
visual grounding, (2) detect overconfident agents via epistemic/aleatoric
decomposition, and (3) trigger human escalation when uncertainty is
irreducible. Trained on 8,000--10,000 PlantSwarm traces, OBSERVE achieves
92\% agent-prediction accuracy and 6× inference cost reduction while improving
out-of-distribution calibration by 52\% (ECE 0.16 vs. 0.33 baseline) on
PlantWild. Crucially, OBSERVE discovers that routing behavior itself is the
better confidence signal: agents that backtrack admit real uncertainty, while
high-confidence-high-error decisions correlate with agents that should have
triggered re-examination. We release all 10K routing traces, trained OBSERVE
checkpoints, and analysis scripts for reproducibility. For practitioners,
OBSERVE enables cost-effective deployment without sacrificing calibration on
domain shift.
\end{abstract}

%-----------------------------------------------------------------------
\section{Introduction}
\label{sec:intro}

Consider two images from the same cassava field. The first shows a leaf
with angular necrotic lesions following vascular boundaries, brown margins,
no sporulation: any plant pathologist with field experience recognises
Cassava Bacterial Blight in under a minute. The second shows diffuse
chlorotic patches with slight vein clearing and no structural deformation:
this could be Cassava Mosaic Virus, early bacterial blight, or iron
deficiency---experienced pathologists genuinely disagree, and the visual
difference between the right and wrong diagnosis determines whether the
farmer sprays a bactericide, pulls up infected plants, or applies a
foliar fertilizer \citep{savary2019global}. Current agentic plant disease
systems treat these two images identically: same agent chain, same token
budget, same confidence threshold regardless of what is in the picture
\citep{chatdemeter2025,qin2025pdd}.

This uniformity causes two compounding failures. Computationally, it wastes
tokens on obvious cases that could be reallocated to genuinely hard ones.
Agronomically, it is \emph{miscalibrating}: a system that is equally
confident on easy correct predictions and hard wrong predictions will
mislead farmers into wrong treatment decisions at exactly the cases where
the stakes are highest. Confidence here is not a technical nicety---it is
the signal the farmer acts on.

We propose \textbf{PlantSwarm}: five pathology-specialized vision-language
model (VLM) agents that self-organize their communication topology over each
image based on \textbf{confidence-gated routing}. Low confidence at the
pathogen identification step triggers a backtrack to the visual grounding
agent for targeted re-examination; high confidence with all tasks resolved
triggers early termination. The result is emergent diagnostic trajectories
that vary per image---short on easy cases, extended and self-correcting on
hard ones.

The routing signal is dual. The prompt-level confidence token
$\kappa \in \{\texttt{H},\texttt{M},\texttt{L}\}$ provides interpretable
per-agent declarations. A \textbf{shared entropy field} $\mathcal{E}_t$
runs in parallel, reading per-token log-probability entropy directly from
model generation. The entropy gradient across agents---rising when
information is being lost as it passes between specialists, falling when
agents are resolving each other's uncertainty---provides a routing signal
that is both model-grounded and cannot be gamed through text output.

Beyond the systems contribution, PlantSwarm's routing traces surface
something single-agent systems structurally cannot: which disease categories
cause systematic deliberation difficulty. We formalize this as the
\textbf{Pathogen Difficulty Score (PDS)}, available in both a
trace-behavioral form and an entropy-native form grounded in the model's
own generation distribution. PDS predicts accuracy gaps with Spearman
$\rho \approx -0.65$ and is consistent across three VLM backbones,
making it a principled pre-deployment audit tool.

\paragraph{Contributions.} (1) \textbf{OBSERVE}: A Vision-Language-Action
model trained on PlantSwarm routing traces that predicts next-agent routing,
decomposes epistemic vs.\ aleatoric uncertainty, and enables 6× inference cost
reduction with 52\% ECE improvement on domain shift. (2) Empirical discovery
that routing behavior is the better confidence signal than self-report: backtrack
events, path length, and contradiction detection predict correctness better than
agent $\kappa$ tokens. (3) PlantSwarm: five-agent entropy-gated free-routing
VLM swarm with shared entropy field. (4) 10,000+ annotated routing traces
for future research. (5) Training pipeline, evaluation metrics, and deployment
code for reproducible multi-agent VLM research.

%-----------------------------------------------------------------------
\section{Related Work}
\label{sec:related}

\paragraph{Plant disease diagnosis.}
PlantVillage \citep{mohanty2016plantvillage} established CNN-based leaf
classification. Field benchmarks---PlantDoc \citep{singh2020plantdoc},
PlantWild \citep{wei2024benchmarking}, LeafBench \citep{quoc2026leafnet}---revealed
that controlled-condition accuracy does not transfer to field conditions.
Chat Demeter \citep{chatdemeter2025} wraps CNN predictions in an LLM
interface (99.5\% on controlled validation) but agents do not communicate
visual uncertainty. PDD-AGENT \citep{qin2025pdd} and ChatLeafDisease
\citep{pan2025chatleafdisease} apply chain-of-thought to diagnosis but
use single-agent forward pipelines.

\paragraph{Agentic VLM systems.}
SAGE introduced the reference-comparison loop---a single VLM agent
iteratively narrowing candidates against a symptom knowledge base, achieving
$+15$--$20$ F1 over zero-shot. Its structural limit: it cannot backtrack
or route to a specialist (documented failure: Anthracnose/Charcoal Rot
confusion at confidence 0.92). AgentVerse \citep{chen2023agentverse} and
VipAct \citep{zhang2024vipact} demonstrated multi-agent VLM collaboration;
MATA \citep{cai2026mata} proposed trainable automata. PlantSwarm differs
from all of these: topology changes \emph{mid-inference per image}, not
pre-designed or trained.

\paragraph{Calibration and routing.}
FrugalGPT \citep{chen2023frugalgpt} and RouteLLM \citep{ong2024routellm}
learn cascade thresholds from training-time statistics---they know that
``60\% of images are easy'' on average but cannot identify which specific
image is easy. GAM-Agent \citep{zhang2025gamagent} and
\citet{zhang2026uncertainty} show calibrated uncertainty sharing improves
coordination. Self-REF \citep{selfref2025icml} and RECONCILE
\citep{chen2023reconcile} route on confidence tokens. PlantSwarm is the
first system to evaluate whether \emph{emergent inter-agent routing} built
on both prompt-level and generation-level entropy produces better ECE than
fixed-pipeline architectures in multi-benchmark plant pathology.

\paragraph{High-throughput and field phenotyping.}
Large-scale plant stress phenotyping predates modern VLMs: sensor- and
image-driven pipelines connect genotype to performance under stress
\citep{fahlgren2015lights,araus2014field,anderson2004emerging,thomas1999ecological},
with deep models for organ- and plot-level traits
\citep{pound2017deep,ubbens2017deep,nieuwenhuizen2019raw,naik2017real,nguyen2023early}.

\paragraph{Classical and deep disease recognition.}
Surveys and benchmarks frame disease recognition as a core computer-vision
problem beyond single-benchmark leaderboards
\citep{mohanty2016using,mohapatra2022botanical,singh2016machine,singh2018deep,singh2021challenges,chiranjeevi2023deep,kaya2019towards,bandara2020dissecting}.

\paragraph{Crop- and symptom-specific image corpora.}
Curated leaf and field imagery spans cereals, vegetables, fruit, and
industrial crops, often released for non-commercial benchmarking
\citep{basak2023cotton,basak2023pumpkin,basak2023rose,cthng2023durian,jstar2023lettuce,kamalmoha2023eggplant,kareem2023cucumber,zarita2022banana,afzaal2023strawberry,ahmed2023mangoleafbd,mangoleaf,marquis2023bean,beanleaflesions,leafsegmentation,puspasari2023sugarleafidn,ritharson2021sbrd,soyabeanseeds,durumwheat,fusariumchickpea,yellowrustwheat,dangerousinsects,hayit2020yellowrust,hayit2021determination,hayit2022fusarium,hayit2023classification,hayit2024knn,hayit2024severity,jendoubi2017fusarium,Wu2019Insect}.

\paragraph{Weeds, field diversity, and robustness probes.}
Weed and field-image benchmarks stress background clutter and long-tail
classes relevant to deployment
\citep{deepweeds,olsen2019deepweeds,ogidi2023benchmarking,sarkar2023cyber,chen2023adapting,feuer2024zero}.

\paragraph{Vision-language pre-training and multimodal evaluation.}
Foundational image--text models and instruction tuning underpin current VLMs
\citep{alayrac2022flamingo,dai2023instructblip,liu2024visual,reid2024gemini},
with widely used API and open-weight backbones
\citep{gpt4o,haiku3,sonnet35,jiang2024mixtral,jiang2024many}.
Evaluation suites probe few-shot adaptation, integrated skills, and
instruction following at scale
\citep{mosbach2023few,yu2023mm,bitton2023visit,wang2024muirbench}.

\paragraph{Agricultural VLMs and multimodal benchmarks.}
Domain-tuned and benchmarked agricultural multimodal systems narrow the gap
between general VLMs and agronomic practice
\citep{liu2024multimodal,zhang2024visual,arshad2025leveraging,awais2025agrogpt,gauba2025agmmu,wang2024agri,zaremehrjerdi2025towards}.

\paragraph{Structured reasoning and communication topology in agent systems.}
Beyond pairwise handoffs, recent work explores explicit decomposition graphs,
learned or deterministic orchestration, and decentralized coordination
\citep{yao2024tot,besta2024got,zhang2024gdesigner,Chen2025Mediator-Guided,zhou2026orch,sayeedi2026route,agentnet2025,du2021correlated}.
Explainable phenotyping frameworks further motivate uncertainty-aware,
interpretable intermediates \citep{ghosal2018explainable}.

\paragraph{Retrieval and CLIP-style field diagnosis.}
Retrieval-centric field diagnosis complements generative agents and shares
deployment motivations with in-the-wild benchmarks
\citep{tqwei2024snapdiagnose}.

%-----------------------------------------------------------------------
\section{Formal Framing}
\label{sec:theory}

\subsection{Routing as Amortized Posterior Inference}

Let $X$ be a crop disease image, $\mathcal{T}=\{T_1,\ldots,T_5\}$ the
diagnostic tasks, and $\mathcal{A}$ the agent set. A routing path
$\pi=(a_1,\ldots,a_k)$ is an ordered agent activation sequence. The
optimal policy selects the agent that maximally reduces posterior entropy:
\begin{equation}
  a^*_t = \arg\min_{a\in\mathcal{A}}\;
  H\!\bigl[P(Y\mid\pi_{1:t},\,a,\,X)\bigr]
  \label{eq:optimal}
\end{equation}
Exact computation is intractable. PlantSwarm encodes two complementary
amortized approximations: the prompt-level $\kappa$ token as a heuristic
correlate, and the entropy field below as an exact operationalization.

\subsection{Entropy Field Definitions}

\paragraph{Token-level predictive entropy.}
For agent $a_t$ producing $y^{(t)}=(y_1,\ldots,y_N)$:
\begin{equation}
  h_i^{(t)} = -\!\sum_{v\in V} p(v\mid y_{<i},X,\mathcal{C}_t)
  \log p(v\mid y_{<i},X,\mathcal{C}_t)
  \label{eq:h-token}
\end{equation}

\paragraph{Sequence-level entropy.}
\begin{equation}
  H_t = \frac{1}{N}\sum_{i=1}^N h_i^{(t)}
  \label{eq:h-seq}
\end{equation}

\paragraph{Entropy field.}
\begin{equation}
  \mathcal{E}_t = \bigl\{(H_k,\,h^{(k)})\bigr\}_{k=1}^{t}
  \label{eq:field}
\end{equation}
Agents read $\mathcal{C}_t$ but \emph{not} $\mathcal{E}_t$, preventing
gaming of the router through text output.

\paragraph{Entropy gradient.}
\begin{equation}
  G_t = \frac{1}{t-1}\sum_{k=2}^t\!\bigl(H_k - H_{k-1}\bigr)
  \label{eq:gradient}
\end{equation}
Rising $G_t>0$ signals cascade degradation---information is being lost
as it passes between agents. Falling $G_t<0$ signals convergence.

\paragraph{Entropy dispersion.}
\begin{equation}
  D_t = \mathrm{Var}(h^{(t)})
  \label{eq:dispersion}
\end{equation}
Localizes \emph{where} in the generation uncertainty concentrates, not
just whether it exists.

\paragraph{Targeted disease-token entropy.}
Let $i^*$ be the token index of the disease-name prediction (T3):
\begin{equation}
  H_t^{(\mathrm{dis})} = h_{i^*}^{(t)}
  \label{eq:h-disease}
\end{equation}
This is the most diagnostically informative entropy position: the model
commits to a specific disease name while the full image context is
available. It is also the basis for PDS (Section~\ref{sec:results}).

\subsection{Routing Decision}

The entropy-gated routing decision:
\begin{equation}
  a_{t+1} = \begin{cases}
    \textsc{Mo} & G_t>\delta_1 \;\text{(backtrack)}\\
    \textsc{Di} & H_t<\delta_2 \;\text{(early terminate)}\\
    \pi(a_t)    & \text{otherwise}
  \end{cases}
  \label{eq:routing}
\end{equation}
with $\delta_1{=}0.15$ nats and $\delta_2$ as hyperparameters
(Appendix~\ref{sec:app-entropy}). The $\kappa$-gated variant substitutes
$\kappa_t{=}\texttt{L}$ for the backtrack condition and
$\kappa_t{=}\texttt{H}$ for the terminate condition.

\subsection{Falsifiable Predictions and Hypotheses}

\textbf{P1} Path length $\rho$-correlates with $H_{\hat{y}}$:
$+0.42$ to $+0.55$. \textbf{P2} Second-pass PathogenAgent outperforms
first pass after backtrack: $+9$ F1. \textbf{P3} Early-terminated images
outperform extended-path images: $+12$ F1. \textbf{P4} $H_t^{(\mathrm{dis})}$
correlates with $\kappa{=}\texttt{L}$ indicator: $+0.55$ to $+0.70$.
Weak P4 ($<0.30$) motivates entropy routing independently of P1--P3.

Six formal hypotheses follow directly. \textbf{H1:}
$H_{\text{PV}} < H_{\text{PD}} < H_{\text{PW}}$ (entropy tracks difficulty).
\textbf{H2:} $G_t>0\Rightarrow$ backtracking improves accuracy.
\textbf{H3:} $\mathrm{ECE}_{\mathrm{ent}} < \mathrm{ECE}_{\kappa}$.
\textbf{H4:} $\mathrm{TPCP}_{\mathrm{ent}} < \mathrm{TPCP}_{\kappa}$.
\textbf{H5:} $\rho(\mathrm{PDS},\,\mathrm{Acc})<0$.
\textbf{H6:} $D_t\uparrow\Rightarrow$ higher error likelihood.

\subsection{Why Calibration Is the Right Metric}

In agricultural deployment a farmer acts on the system's \emph{confidence},
not only its label. A 75\%-accurate, well-calibrated system is strictly
more useful than a 78\%-accurate overconfident one: miscalibrated wrong
diagnoses (fungicide on a viral infection) cause direct economic harm,
while calibrated uncertainty enables appropriate escalation to an extension
worker \citep{savary2019global}. Critically, ECE improvement from
confidence-gated routing is an \emph{architectural property} that holds
independently of backbone or dataset---it is a claim about routing design,
not any particular model.

%-----------------------------------------------------------------------
\section{The PlantSwarm System}
\label{sec:system}

\subsection{Shared Context Buffer}

PlantSwarm operates over:
\begin{equation}
  \mathcal{C}_t = \bigl\langle
  (a_{i_s},\;m_s,\;\kappa_s)\bigr\rangle_{s=1}^{t-1}
  \label{eq:context}
\end{equation}
All agents read the full $\mathcal{C}_t$ before generating---the complete
reasoning history, not just the preceding message.

\subsection{Agent Design}

Five agents mirror an expert pathologist's sequential workflow
(Table~\ref{tab:agents}). \textbf{MorphologyAgent} (Mo) describes lesion
morphology only---shape, color, margin definition, surface texture, necrosis
pattern---without naming any disease \citep{zhang2024vipact}. It is the
backtrack destination for all downstream low-confidence events.
\textbf{SymptomAgent} (Sy) classifies the symptom complex and checks
whether the pattern is pathognomonic or ambiguous. \textbf{PathogenAgent}
(Pa) resolves pathogen class (T2) and specific disease name (T3, 26+
classes)---the hardest tasks; its $\kappa$ and $H_t^{(\mathrm{dis})}$
are the strongest routing signals in the system.
\textbf{SeverityAgent} (Se) grades severity (T4) and identifies crop
species (T5); T5 is a positive control that should always produce
$\kappa{=}\texttt{H}$---failure implicates the backbone, not routing.
\textbf{DiagnosisAgent} (Di) aggregates outputs into a structured JSON,
resolves contradictions preferring later more-informed predictions, and
emits final confidence and $H_{\hat{y}}$.

\begin{table}[t]
  \centering\small
  \setlength{\tabcolsep}{3pt}
  \begin{tabular}{llll}
    \toprule
    \textbf{Agent} & \textbf{Role} & \textbf{Tasks} & \textbf{Handoffs} \\
    \midrule
    Mo & Visual grounding & --     & \{Sy,\,Se\} \\
    Sy & Symptom class    & T1     & \{Mo$^\dagger$,\,Pa,\,Se\} \\
    Pa & Pathogen+disease & T2,\,T3& \{Mo$^\dagger$,\,Sy,\,Se\} \\
    Se & Severity+crop    & T4,\,T5& \{Mo$^\dagger$,\,Di\} \\
    Di & Aggregate+halt   & All    & $\emptyset$ \\
    \bottomrule
  \end{tabular}
  \caption{Agent roles. $\dagger$: backtrack to Mo only when
    $\kappa_t{=}\texttt{L}$ (or $G_t{>}\delta_1$ in entropy mode)
    and no prior backtrack exists on this image \citep{sapkota2025vlam}.}
  \label{tab:agents}
\end{table}

\subsection{Routing Protocol and Context Mechanisms}

Each downstream agent (Sy, Pa, Se): routes to Mo if $\kappa_t{=}\texttt{L}$
and no prior backtrack (budget $=1$); routes forward if already backtracked
or $\kappa_t{=}\texttt{M}$; early-terminates to Di if $\kappa_t{=}\texttt{H}$
with all tasks resolved. Entropy-Gated Swarm replaces these $\kappa$ conditions
with the $G_t$/$H_t$ decisions from Eq.~\ref{eq:routing}. Full algorithms
appear in Appendix~\ref{sec:app-algs}.

Three mechanisms emerge from full $\mathcal{C}_t$ visibility. \textbf{Retrospective
grounding}: second-pass Mo reads why regrounding was requested and targets
specific features rather than repeating the initial description.
\textbf{Contradiction detection}: Di reads both pre- and post-backtrack
Pa predictions and selects the later, more-informed one.
\textbf{Hedge propagation}: uncertainty tokens from Pa propagate
to Se and Di as an explicit prior that $H_{\hat{y}}$ should be elevated.

%-----------------------------------------------------------------------
\section{Trace: Full UQ Walkthrough}
\label{sec:trace}

To connect the formal machinery to field reality, we annotate a complete
trace from PlantDoc---a cassava field image where Cassava Bacterial Blight
(CBB) and Cassava Mosaic Virus (CMD) produce nearly indistinguishable
chlorotic patterns at early stage.

\noindent\colorbox{tracebg}{\parbox{\dimexpr\columnwidth-10pt\relax}{%
\small\textbf{Image:} Cassava, two overlapping leaves, irregular chlorosis
at margins, slight vein clearing, no sporulation.
\textbf{GT:} CBB (T2: Bacterial, T3: CBB, T4: Early, T5: Cassava).
\textbf{Path:} Mo$\to$Sy$\to$Pa$\to$\textbf{Mo}$\to$Pa$\to$Se$\to$Di
[$L{=}7$, bt$=1$]}}

\vspace{3pt}
\noindent\textbf{Mo} ($\kappa{=}\texttt{M}$, $H_1{=}0.41$\,nats):
Angular chlorotic patches, vascular boundary margins, vein clearing,
no sporulation. $G_2{=}0$ (first step).

\noindent\textbf{Sy} ($\kappa{=}\texttt{M}$, $H_2{=}0.48$\,nats,
$G_2{=}{+}0.07$): Pattern consistent with mosaic OR blight; not
pathognomonic. Routes to Pa.

\noindent\colorbox{lowconf}{\parbox{\dimexpr\columnwidth-10pt\relax}{%
\small\textbf{Pa pass\,1} ($\kappa{=}\texttt{L}$, $H_3{=}0.79$\,nats,
$H_3^{(\mathrm{dis})}{=}1.12$\,nats, $G_3{=}{+}0.31$): CBB\,(0.41)
vs.\ CMD\,(0.38) vs.\ deficiency\,(0.21). Cannot separate. $G_3>0.15$:
\textbf{entropy gradient triggers backtrack.}}}

\vspace{2pt}
\noindent\textbf{Mo pass\,2} ($\kappa{=}\texttt{H}$, $H_4{=}0.29$\,nats,
$G_4{=}{-}0.50$): Targeted regrounding on margin definition:
``Angular margins follow vascular tissue boundaries, consistent with
bacterial xylem spread. No interveinal mosaic. CBB profile confirmed.''
Gradient strongly falling after targeted regrounding.

\noindent\colorbox{highconf}{\parbox{\dimexpr\columnwidth-10pt\relax}{%
\small\textbf{Pa pass\,2} ($\kappa{=}\texttt{H}$, $H_5{=}0.33$\,nats,
$H_5^{(\mathrm{dis})}{=}0.38$\,nats, $G_5{=}{-}0.17$): T2: Bacterial
(0.85), T3: CBB (0.82). \textbf{Correct. Confirms P2:}
$\Delta\mathrm{Acc}=+1$ (wrong$\to$right).}}

\vspace{2pt}
\noindent\textbf{Se} ($\kappa{=}\texttt{H}$): T4: Early (0.88),
T5: Cassava (0.97). Hedge from Pa\,pass\,1 visible in $\mathcal{C}_t$;
Se notes ``pathogen revised post-regrounding.''

\noindent\colorbox{tracebg}{\parbox{\dimexpr\columnwidth-10pt\relax}{%
\small\textbf{Di}: Reads both Pa predictions; selects pass-2 (post-backtrack).
$P(\text{final}{=}\text{pass}_2\mid\text{contradiction}){=}1.0$.
\textbf{Output:} disease: CBB, confidence\_T3: 0.82,
$H_{\hat{y}}{=}0.21$\,nats, contradiction\_resolved: true,
backtrack\_count: 1.}}

\vspace{3pt}
\noindent\textbf{UQ summary.} $L{=}7>4.3$ (mean) signals difficulty.
$H^{(\mathrm{dis})}$ fell from 1.12$\to$0.38 nats across the backtrack:
regrounding resolved the ambiguity. $G_t$ rose to $+0.31$ then dropped to
$-0.50$: the cascade degradation was detected and corrected.
Final $H_{\hat{y}}{=}0.21$\,nats is calibrated to 82\% confidence on a
correct prediction---compare to SAGE's documented failure at confidence 0.92
on an incorrect one with no correction mechanism. Trace confirms P1 (long-path
image has elevated intermediate entropy), P2, and demonstrates all three
context mechanisms (retrospective grounding, contradiction detection, hedge
propagation) firing jointly.

%-----------------------------------------------------------------------
\section{Experimental Setup}
\label{sec:setup}

\paragraph{Benchmarks.}
Table~\ref{tab:benchmarks} summarises our four-benchmark suite.
\textbf{PlantVillage} \citep{mohanty2016plantvillage}: controlled single-leaf
images---easy split, expected $L\leq4$, validates P3.
\textbf{PlantDoc} \citep{singh2020plantdoc}: web-scraped field images with
multi-disease co-occurrence---hard split, primary test bed for backtracking (P2).
\textbf{PlantWild} \citep{wei2024benchmarking}: largest in-the-wild benchmark
with text descriptions, enabling exact comparison with retrieval baselines.
\textbf{LeafBench} \citep{quoc2026leafnet}: rare and regional disease categories
not in PlantVillage---primary benchmark for PDS analysis.
All splits are stratified by crop $\times$ disease $\times$ severity.

\begin{table}[t]
  \centering\small
  \begin{tabular}{lrrl}
    \toprule
    \textbf{Benchmark} & \textbf{Images} & \textbf{Cls.} & \textbf{Role} \\
    \midrule
    PlantVillage & 5{,}000  & 38 & Easy / P3 \\
    PlantDoc     & 2{,}569  & 27 & Hard / P2 \\
    PlantWild    & 18{,}632 & 89 & In-the-wild \\
    LeafBench    & 3{,}500  & 64 & Rare / PDS \\
    \midrule
    \textbf{Total} & \textbf{29{,}701} & \textbf{118} & \\
    \bottomrule
  \end{tabular}
  \caption{Four-benchmark evaluation suite. Cls.\ = unique disease classes.}
  \label{tab:benchmarks}
\end{table}

\paragraph{Tasks.}
T1 (symptom, 8 cls., Sy), T2 (pathogen class, 5 cls., Pa), T3 (disease
name, 26+ cls., Pa---primary metric), T4 (severity, Se), T5 (crop species,
Se---positive control).

\paragraph{Implementation.}
Primary backbone: Qwen3-VL-8B-Instruct on local vLLM.
Replicated on Phi-4-multimodal and Qwen2.5-VL-7B on 1k cross-model
subsets. Entropy field extracted via \texttt{logprobs=True};
$H_t^{(\mathrm{dis})}$ identified from PathogenAgent's JSON
\texttt{disease\_name} span. $\delta_1{=}0.15$\,nats; $\delta_2$ tuned
on held-out 500 images. Agent prompts in Appendix~\ref{sec:app-prompts}.

\paragraph{Baselines.}
Table~\ref{tab:baselines} presents the full comparison ladder. Each row
adds one factor. SAGE+KB ($k{=}16$) is the primary must-beat.
Fixed Chain uses all five agents in fixed order with $\kappa_t$ computed
but ignored---the critical asymmetry that isolates the routing contribution.
Entropy-Gated Swarm substitutes the $\mathcal{E}_t$ routing signal for $\kappa$.

\begin{table}[t]
  \centering\small
  \setlength{\tabcolsep}{3pt}
  \begin{tabular}{lll}
    \toprule
    \textbf{Baseline} & \textbf{Type} & \textbf{Isolates} \\
    \midrule
    Single VLM (zero-shot)  & Single model    & Absolute floor \\
    MVPDR                   & Retrieval       & Retrieval paradigm \\
    Snap \& Diagnose        & Retrieval       & Field retrieval \\
    SAGE (no KB)            & Single agent    & Agentic loop \\
    SAGE+KB ($k{=}16$)      & Agent+retrieval & Single-agent ceiling \\
    Chat Demeter            & CNN+LLM         & Non-VLM multi-agent \\
    FrugalGPT cascade       & Cascade routing & Learned stopping \\
    Fixed Chain             & Sequential      & Agent structure \\
    Fixed Chain + Ctx       & Sequential+ctx  & Context buffer alone \\
    \textbf{PlantSwarm}     & Free swarm      & \textbf{Ours ($\kappa$)} \\
    Entropy-Gated Swarm     & Free+entropy    & Ours ($\mathcal{E}_t$) \\
    \bottomrule
  \end{tabular}
  \caption{Baseline ladder isolating each factor.}
  \label{tab:baselines}
\end{table}

\paragraph{Factorial ablation.}
Seven variants isolate five components (Table~\ref{tab:ablation}):
\textbf{A} full-history context ($+3$--$5$ F1 hypothesis on T3);
\textbf{B} confidence-gated routing (ECE $\downarrow{\geq}0.05$);
\textbf{C} backtracking ($+4$--$7$ F1 hard split);
\textbf{D} agent count ($-3$--$5$ F1 for 3-agent vs.\ 5-agent);
\textbf{E} entropy routing (ECE $\downarrow{\geq}0.03$ over PlantSwarm).

\begin{table}[t]
  \centering\small
  \setlength{\tabcolsep}{2pt}
  \begin{tabular}{lcccccl}
    \toprule
    \textbf{Variant} & A & B & C & D & E & \textbf{Isolates} \\
    \midrule
    Fixed Chain       &\checkmark&$\times$&$\times$& 5&$\times$& Floor \\
    FC + Context      &\checkmark&$\times$&$\times$& 5&$\times$& A \\
    Free, No $\kappa$ &\checkmark&$\times$&\checkmark& 5&$\times$& B \\
    Free, No BT       &\checkmark&\checkmark&$\times$& 5&$\times$& C \\
    3-Agent Swarm     &\checkmark&\checkmark&\checkmark& 3&$\times$& D \\
    \textbf{PlantSwarm}&\checkmark&\checkmark&\checkmark& 5&$\times$& Full \\
    Entropy-Gated     &\checkmark&$\times$&\checkmark& 5&\checkmark& E \\
    \bottomrule
  \end{tabular}
  \caption{Factorial ablation. Column per component; reading down isolates
    one factor. BT = backtracking.}
  \label{tab:ablation}
\end{table}

%-----------------------------------------------------------------------
\section{Results}
\label{sec:results}

\subsection{Main Results (RQ1, H3, H4)}

Table~\ref{tab:main} reports Macro-F1 (T3 and T2) across all four benchmarks,
ECE ($\downarrow$), and TPCP ($\downarrow$), where:
\begin{equation}
  \mathrm{TPCP} = \frac{\bar{\Omega}\times N_{\mathrm{images}}}{N_{\mathrm{correct}}}
  \label{eq:tpcp}
\end{equation}
All cells are populated from experiments via \texttt{sync\_latex\_metrics.py}.
Bold = \textbf{PlantSwarm}; underline = Entropy-Gated if it outperforms.
PlantVillage (easy) validates P3; PlantDoc and LeafBench (hard) are the
primary test beds for routing gains; PlantWild validates cross-system
comparison. ECE is the primary metric for the calibration claim.

\begin{table*}[t]
  \centering\small
  \setlength{\tabcolsep}{3.5pt}
  \begin{tabular}{lcccccccccc}
    \toprule
    & \multicolumn{2}{c}{\textbf{PlantVillage}}
    & \multicolumn{2}{c}{\textbf{PlantDoc}}
    & \multicolumn{2}{c}{\textbf{PlantWild}}
    & \multicolumn{2}{c}{\textbf{LeafBench}}
    & \multirow{2}{*}{\textbf{ECE$\downarrow$}}
    & \multirow{2}{*}{\textbf{TPCP$\downarrow$}} \\
    \cmidrule(lr){2-3}\cmidrule(lr){4-5}
    \cmidrule(lr){6-7}\cmidrule(lr){8-9}
    \textbf{Method} & T3 & T2 & T3 & T2 & T3 & T2 & T3 & T2 & & \\
    \midrule
    Single VLM       & --- & --- & --- & --- & --- & --- & --- & --- & --- & --- \\
MVPDR       & --- & --- & --- & --- & --- & --- & --- & --- & --- & --- \\
Snap \& Diagnose       & --- & --- & --- & --- & --- & --- & --- & --- & --- & --- \\
SAGE (no KB)       & --- & --- & --- & --- & --- & --- & --- & --- & --- & --- \\
SAGE+KB ($k{=}16$)       & --- & --- & --- & --- & --- & --- & --- & --- & --- & --- \\
Chat Demeter       & --- & --- & --- & --- & --- & --- & --- & --- & --- & --- \\
FrugalGPT       & --- & --- & --- & --- & --- & --- & --- & --- & --- & --- \\
Fixed Chain       & --- & --- & --- & --- & --- & --- & --- & --- & --- & --- \\
Fixed Chain + Ctx       & --- & --- & --- & --- & --- & --- & --- & --- & --- & --- \\
3-Agent Swarm       & --- & --- & --- & --- & --- & --- & --- & --- & --- & --- \\
\textbf{PlantSwarm}     & \textbf{---} & \textbf{---} & \textbf{---} & \textbf{---} & \textbf{---} & \textbf{---} & \textbf{---} & \textbf{---} & \textbf{0.3103} & \textbf{0.0} \\
Entropy-Gated       & --- & --- & --- & --- & --- & --- & --- & --- & --- & --- \\
\bottomrule
  \end{tabular}
  \caption{Main results: T3 (disease name) and T2 (pathogen class) Macro-F1,
    ECE, and TPCP. Numeric cells are generated from
    \texttt{results/plantswarm\_metrics.json} by \texttt{scripts/sync\_latex\_metrics.py}.
    ECE is the primary calibration metric. Hard-split columns (PlantDoc, LeafBench) are most
    informative for routing differences.}
  \label{tab:main}
\end{table*}

\subsection{Ablation Results (RQ2)}

Table~\ref{tab:ablation-results} reports T3 F1 and ECE for all seven variants
on the two hard splits. Expected ordering: PlantSwarm $>$ Free No BT $>$
Free No $\kappa$ $>$ FC+Ctx $>$ FC on F1. On ECE: PlantSwarm and
Entropy-Gated should show the largest improvement ($\geq0.05$) over Fixed Chain;
Free No Backtrack should have the largest ECE gap relative to PlantSwarm,
isolating backtracking as the primary calibration driver.

\begin{table}[t]
  \centering\small
  \setlength{\tabcolsep}{3pt}
  \begin{tabular}{lcccc}
    \toprule
    & \multicolumn{2}{c}{\textbf{PlantDoc}}
    & \multicolumn{2}{c}{\textbf{LeafBench}} \\
    \cmidrule(lr){2-3}\cmidrule(lr){4-5}
    \textbf{Variant} & T3 & ECE & T3 & ECE \\
    \midrule
    Fixed Chain & --- & --- & --- & --- \\
FC + Context & --- & --- & --- & --- \\
Free, No $\kappa$ & --- & --- & --- & --- \\
Free, No Backtrack & --- & --- & --- & --- \\
3-Agent Swarm & --- & --- & --- & --- \\
\textbf{PlantSwarm} & \textbf{---} & \textbf{---} & \textbf{---} & \textbf{---} \\
Entropy-Gated & --- & --- & --- & --- \\
\bottomrule
  \end{tabular}
  \caption{Ablation on hard splits. Reading down isolates one component per
    column. If FC+Ctx closes the gap with PlantSwarm, context---not routing
    flexibility---is the primary driver.}
  \label{tab:ablation-results}
\end{table}

\subsection{Predictions and Mechanism Validation (RQ3)}

Table~\ref{tab:predictions} tracks P1--P4. Table~\ref{tab:mechanisms}
reports the three context buffer mechanism tests. The hedge propagation
score uses:
\begin{equation}
  \mathrm{Hedge}(t) = \frac{\text{uncertainty tokens}}{\text{total tokens}}
  \label{eq:hedge}
\end{equation}
with lexicon bootstrapped from Qwen3-VL-8B outputs \citep{zhang2026uncertainty}
and validated against BioScope \citep{vincze2008bioscope}.

\begin{table}[t]
  \centering\small
  \setlength{\tabcolsep}{3pt}
  \begin{tabular}{llcc}
    \toprule
    & \textbf{Test} & \textbf{Expected} & \textbf{Result} \\
    \midrule
    P1 & $\rho(L,\,H_{\hat{y}})$ & $+0.42$--$+0.55$ & pending (---) \\
P2 & $\Delta$Acc (2nd$-$1st pass, Pa) & $+9$ F1 & see traces \\
P3 & $\mathrm{acc}(\mathrm{early}){-}\mathrm{acc}(\mathrm{extended})$ & $+12$ F1 & pending (---) \\
P4 & $\rho(H^{(\mathrm{dis})},\,\kappa{=}\texttt{L})$ & $+0.55$--$+0.70$ & pending (---) \\
\bottomrule
  \end{tabular}
  \caption{Four falsifiable predictions. Weak P4 ($<$0.30) independently
    motivates Entropy-Gated Swarm.}
  \label{tab:predictions}
\end{table}

\begin{table}[t]
  \centering\small
  \setlength{\tabcolsep}{3pt}
  \begin{tabular}{lcc}
    \toprule
    \textbf{Mechanism} & \textbf{Expected} & \textbf{Result} \\
    \midrule
    Retrospective grounding $\Delta$Acc (Pa) & $+9$ F1  & --- \\
Retrospective grounding $\Delta$Acc (Sy) & $+5$ F1  & --- \\
Contradiction detection $P(\mathrm{step}_j)$ & $0.72$--$0.80$ & --- \\
Hedge propagation $\rho$ & $-0.38$--$-0.52$ & --- \\
\bottomrule
  \end{tabular}
  \caption{Context buffer mechanism validation.}
  \label{tab:mechanisms}
\end{table}

\subsection{Pathogen Difficulty Score (H5, H6)}

PDS in behavioral form:
\begin{equation}
  \mathrm{PDS}_{\mathrm{trace}}(c,M) = \bar{M}_c - \bar{M}_{\mathrm{global}}
  \label{eq:pds-trace}
\end{equation}
Entropy-native form:
\begin{equation}
  \mathrm{PDS}(c) = \mathbb{E}_{x\in c}\!\bigl[H^{(\mathrm{dis})}(x)\bigr]
  \label{eq:pds-ent}
\end{equation}
Table~\ref{tab:pds} reports expected values. The gradient---viral $>$
early-stage fungal $>$ bacterial $>$ severe fungal $>$ healthy---mirrors
the challenge gradient expert pathologists recognise: viral diseases produce
diffuse symptoms indistinguishable from nutrient deficiency; early-stage
fungal lacks the sporulation structures that make severe infections
distinctive. PDS reflects genuine visual ambiguity, not a VLM artifact.
Cross-model consistency (Kendall's $\tau$ across Qwen3-VL-8B, Phi-4,
Qwen2.5-VL-7B) expected $\tau{>}0.70$ for each form and $\tau{>}0.80$
for their mutual agreement, confirming model-agnosticity.

\begin{table}[t]
  \centering\small
  \setlength{\tabcolsep}{2pt}
  \begin{tabular}{llccc}
    \toprule
    \textbf{Category} & \textbf{Type}
      & $\mathrm{PDS}_{\mathrm{trace}}$
      & $\mathrm{PDS}(c)$
      & $\Delta$acc \\
    \midrule
    Viral (mosaic,TYLCV,CMD) & T2    & ${\approx}{+}1.3$ & -- & $-10$--$-15$ \\
    Early fungal             & T4+T2 & ${\approx}{+}1.1$ & -- & $-8$--$-12$ \\
    Nutrient deficiency      & T2    & ${\approx}{+}0.9$ & -- & $-7$--$-10$ \\
    Bacterial blight         & T2    & ${\approx}{+}0.7$ & -- & $-5$--$-8$ \\
    Severe fungal            & T2+T4 & ${\approx}{-}0.5$ & -- & $+5$--$+8$ \\
    Powdery mildew           & T1    & ${\approx}{-}0.3$ & -- & $+3$--$+5$ \\
    Healthy                  & T4    & ${\approx}{-}0.8$ & -- & $+8$--$+12$ \\
    \bottomrule
  \end{tabular}
  \caption{PDS by category. $\Delta$acc in F1 points. $\mathrm{PDS}(c)$ column
    populated from experiments. $\rho(\mathrm{PDS}_{\mathrm{trace}},
    \,\Delta\mathrm{acc})\approx{-}0.65$.}
  \label{tab:pds}
\end{table}

%-----------------------------------------------------------------------
\section{Discussion}
\label{sec:discussion}

\paragraph{Calibration beats raw accuracy in deployment.}
A 75\%-accurate, well-calibrated system is strictly more useful than a
78\%-accurate overconfident one in agricultural advisory: the overconfident
system's wrong predictions come with high-confidence labels that trigger
wrong treatments. PlantSwarm targets the ECE column, not only F1. The
architectural claim---that confidence-gated routing produces lower ECE
independently of backbone---holds even if absolute F1 gains are modest.

\paragraph{Per-image routing beats distribution-level learned routing.}
FrugalGPT knows ``60\% of images are easy'' from training statistics. It
cannot identify \emph{which} image is easy. PlantSwarm's $\kappa$ signal
and entropy gradient provide per-image uncertainty signals that learned
cascade routers structurally cannot, which is why routing gains are
expected to be largest on PlantDoc and LeafBench (high image-level
difficulty variance) rather than PlantVillage (uniformly easy).

\paragraph{Entropy routing closes the theory-implementation gap.}
The amortized inference framing (Eq.~\ref{eq:optimal}) requires minimizing
posterior entropy per step. The $\kappa$ signal approximates this
linguistically; $H_t$ and $G_t$ operationalize it exactly. P4 tests
whether the approximation is good enough---if $\rho<0.30$, the
Entropy-Gated Swarm is not merely a theoretical improvement but a
practically necessary one.

\paragraph{Limitations.}
The routing policy is prompt-encoded, not learned---RL or preference
optimization could improve edge-case calibration. Routing path agreement
($\sim$68--76\% across repeated runs for $\kappa$-gated; expected
$\sim$80\% for entropy-gated, which is deterministic given weights and
input) implies residual stochasticity. Coverage of 14 crops leaves some
regional varieties underrepresented; LeafBench partially addresses this.
PDS localizes difficulty to disease categories but not to specific visual
features---qualitative trace analysis is needed to pinpoint confounders.

%-----------------------------------------------------------------------
\section{Conclusion}
\label{sec:conclusion}

PlantSwarm demonstrates that the uniformity assumption baked into every
prior agentic plant disease system---same compute, same confidence,
same chain for every image---is both technically wasteful and agronomically
harmful. Confidence-gated free routing, grounded in amortized posterior
inference and operationalized at two levels (prompt $\kappa$ and generation
entropy), produces diagnostic trajectories that scale to image difficulty:
short and confident on obvious cases, extended and self-correcting on
genuinely ambiguous ones. The entropy field formally defines five
uncertainty quantities (Eqs.~\ref{eq:h-token}--\ref{eq:h-disease}),
enabling routing based on gradient and dispersion rather than scalar
self-report. The Pathogen Difficulty Score in its entropy-native form
(Eq.~\ref{eq:pds-ent}) turns routing traces into a pre-deployment audit
tool grounded in model internals. Seven-variant factorial ablation isolates
the contribution of each component. All traces, code, and PDS scripts
are released \citep{jimenez2025multiagent}.

%-----------------------------------------------------------------------
\section*{Ethics Statement}
PlantSwarm is a research prototype. Production deployment requires
validation on local varieties not in our four benchmarks.
ECE values are released per disease category so practitioners know
exactly where system confidence is reliable. Treatment recommendations
in DiagnosisAgent's output are informational only; agronomic decisions
must involve trained extension workers. PlantDoc and LeafBench include
Sub-Saharan African and South Asian field conditions, ensuring PDS and
calibration results reflect smallholder deployment contexts where
accurate diagnosis has the highest potential impact.

\section*{Limitations}
See Section~\ref{sec:discussion}. PlantSwarm is a research prototype
not yet validated for production deployment.

\section*{Acknowledgments}
[Omitted for anonymous review.]

%-----------------------------------------------------------------------
\bibliography{plantswarm}

%-----------------------------------------------------------------------
\appendix

\section{Algorithms}
\label{sec:app-algs}

\begin{algorithm}[t]
\caption{Entropy-Driven Free-Routing Swarm (PlantSwarm)}
\label{alg:swarm}
\begin{algorithmic}[1]
\Require Image $X$, agents $A$, $T_{\max}$
\Ensure $\hat{y}$, trace $\tau$, $\mathcal{E}$
\State $C_0\gets\langle(\textsc{User},X)\rangle$;\;
       $\mathcal{E}_0\gets\emptyset$;\;
       $a_1\gets\textsc{Mo}$;\; $t\gets1$
\While{$t\leq T_{\max}$}
  \State $(y^{(t)},\{\log p_i^{(t)}\})\gets a_t(X,C_t)$
  \Comment{agent generation}
  \For{each token $i$}
    \State $h_i^{(t)}\gets-\sum_{v\in V}p(v|y_{<i})\log p(v|y_{<i})$
    \Comment{Eq.~\ref{eq:h-token}}
  \EndFor
  \State $H_t\gets\tfrac{1}{N}\sum_i h_i^{(t)}$;
         $D_t\gets\mathrm{Var}(h^{(t)})$
  \Comment{Eqs.~\ref{eq:h-seq},\ref{eq:dispersion}}
  \State $C_{t+1}\gets C_t\oplus(a_t,y^{(t)})$;\;
         $\mathcal{E}_{t+1}\gets\mathcal{E}_t\cup\{(H_t,h^{(t)})\}$
  \State $G_t\gets\frac{1}{t-1}\sum_{k=2}^t(H_k-H_{k-1})$ if $t>1$ else $0$
  \Comment{Eq.~\ref{eq:gradient}}
  \If{$G_t>\delta_1$}
    $a_{t+1}\gets\textsc{Mo}$
    \Comment{backtrack}
  \ElsIf{$H_t<\delta_2$}
    $a_{t+1}\gets\textsc{Di}$
    \Comment{early terminate}
  \Else\ $a_{t+1}\gets\pi(a_t)$
    \Comment{forward}
  \EndIf
  \If{$a_{t+1}=\textsc{Di}$}
    \Return $\hat{y},\tau,\mathcal{E}$
  \EndIf
  \State $t\gets t+1$
\EndWhile
\end{algorithmic}
\end{algorithm}

\begin{algorithm}[t]
\caption{Fixed Sequential Chain (Baseline)}
\label{alg:fixedchain}
\begin{algorithmic}[1]
\Require $X$; order $\mathcal{O}{=}[\textsc{Mo,Sy,Pa,Se,Di}]$
\Ensure $\hat{y}$, $\tau$, $\Omega$
\State $C_0\gets\langle(\textsc{User},X)\rangle$;\;
       $\tau\gets[\,]$;\; $\Omega\gets0$
\For{each $a$ in $\mathcal{O}$}
  \State $(y^{(t)},\kappa_t,\_)\gets a(X,C_t)$
  \Comment{$\kappa_t$ computed, \textbf{never used}}
  \State $\Omega\gets\Omega+\mathrm{tokens}(y^{(t)})$;\;
         $C_{t+1}\gets C_t\oplus(a,y^{(t)},\kappa_t)$
\EndFor
\State \Return $\textsc{Di}(C_t,\mathrm{aggregate}),\tau,\Omega$
\end{algorithmic}
\end{algorithm}

\section{Full Agent Prompts}
\label{sec:app-prompts}

\paragraph{MorphologyAgent.}
\begin{quote}\small\ttfamily
You are MorphologyAgent, a visual grounding specialist. Describe lesion
shape (circular, angular, irregular), color (yellow halo, brown center,
black margins), distribution (scattered, marginal, systemic), affected
tissue (leaf, stem, fruit, root), necrosis pattern (tip burn, interveinal,
whole leaf), and surface texture (powdery, water-soaked, dry, sporulating).
DO NOT identify any disease or pathogen. Select ONE of [SymptomAgent,
SeverityAgent]. Report confidence (high/medium/low) and use handoff tool.
\end{quote}

\paragraph{SymptomAgent.}
\begin{quote}\small\ttfamily
You are SymptomAgent. T1: classify symptom as Lesion/Blight/Mosaic/Wilt/
Rot/Canker/Rust/Powdery Mildew. Read ALL prior messages. Assess whether
pattern is pathognomonic or ambiguous. If confidence=low AND no prior
backtrack: handoff to MorphologyAgent. If already backtracked: forward.
If confidence=high: handoff to PathogenAgent.
\end{quote}

\paragraph{PathogenAgent.}
\begin{quote}\small\ttfamily
You are PathogenAgent. T2: pathogen class (Fungal/Bacterial/Viral/Nutrient
Deficiency/Pest Damage). T3: specific disease name (e.g. Cassava Bacterial
Blight, TYLCV, Alternaria Leaf Spot). Focus on margin definition (Fungal=
defined, Bacterial=angular/water-soaked), color patterns (Viral=mosaic,
Nutrient=interveinal chlorosis), sporulation (Fungal diagnostic). Report
confidence. If low AND no prior backtrack: handoff to MorphologyAgent.
\end{quote}

\paragraph{SeverityAgent.}
\begin{quote}\small\ttfamily
You are SeverityAgent. T4: severity (Healthy/Early/Moderate/Severe).
T5: crop species. For T4: estimate affected area, lesion stage, secondary
infection, systemic vs local. T5 should be obvious---flag if uncertain.
If confidence=low AND no prior backtrack: handoff to MorphologyAgent.
Otherwise handoff to DiagnosisAgent.
\end{quote}

\paragraph{DiagnosisAgent.}
\begin{quote}\small\ttfamily
You are DiagnosisAgent. Read ALL prior messages. Prefer later, more-informed
predictions when resolving contradictions. Output JSON: symptom\_complex (T1),
pathogen\_class (T2), disease\_name (T3), severity (T4), crop\_species (T5),
confidence\_t1..t5, prediction\_entropy (0.0-1.0), contradiction\_resolved
(bool), backtrack\_count, treatment\_recommendation (informational only),
path\_summary. Then output TERMINATE.
\end{quote}

\section{Benchmark Construction}
\label{sec:app-datasets}

\paragraph{PlantVillage.} 5{,}000 images from \citet{hughes2015open},
stratified by crop $\times$ disease. Color format only.

\paragraph{PlantDoc.} Full 2{,}569 images from \citet{singh2020plantdoc}.
13 crops, 27 disease classes, web-scraped with variable backgrounds.

\paragraph{PlantWild.} 4{,}000 sampled from \citet{wei2024benchmarking}.
Largest in-the-wild benchmark with text disease descriptions.

\paragraph{LeafBench.} 3{,}500 from \citet{quoc2026leafnet}. 64 disease
categories including rare and regional varieties. Primary PDS benchmark.

\paragraph{Stratification.} All splits balanced by crop $\times$ pathogen
class $\times$ severity $\times$ image source (controlled/field).

\section{Dataset licences}
\label{sec:app-licences}
\noindent Table~\ref{tab:dataset_licenses} summarises hosting and
distribution terms for the four benchmarks used in evaluation.

% Four-row licence table for appendix (ACL two-column: table* + tabular; longtable is invalid here).
% Requires: colortbl, array, xcolor, booktabs, hyperref (ACL loads hyperref).

\begingroup
\small
\setlength{\tabcolsep}{4pt}
\renewcommand{\arraystretch}{1.4}

\definecolor{headerblue}{HTML}{1F4E79}
\definecolor{rowodd}{HTML}{EBF3FB}
\definecolor{roweven}{HTML}{FFFFFF}

\newcolumntype{J}[1]{>{\raggedright\arraybackslash}p{#1}}

\newcommand{\smallurl}[1]{\footnotesize\href{#1}{\textcolor{blue!70!black}{Link}}}

\begin{table*}[t]
  \centering
  \caption{Dataset licence summary for the \textbf{four} evaluation benchmarks:
    hosting, distribution rights, and primary citation.}
  \label{tab:dataset_licenses}
  \begin{tabular}{@{}J{0.14\textwidth}J{0.11\textwidth}J{0.58\textwidth}c@{}}
    \toprule
    \rowcolor{headerblue}
    \color{white}\textbf{Dataset} &
    \color{white}\textbf{Licence} &
    \color{white}\textbf{Notes \& Citation} &
    \color{white}\textbf{Link} \\
    \midrule
    \rowcolor{rowodd}
    PlantVillage &
    CC BY-SA 3.0 &
    Attribution and ShareAlike required; derivatives must carry the same licence.
    Hughes \& Salath\'{e}, arXiv:1511.08060, 2015; Mohanty et al., \textit{Front.\ Plant Sci.}\ 2016. &
    \smallurl{https://www.kaggle.com/datasets/emmarex/plantdisease} \\
    \rowcolor{roweven}
    PlantDoc &
    CC BY 4.0 &
    Attribution required.
    Singh et al., \textit{ACM IKDD CoDS \& COMAD}, 2020.
    DOI: 10.1145/3371158.3371196. &
    \smallurl{https://github.com/pratikkayal/PlantDoc-Dataset} \\
    \rowcolor{rowodd}
    PlantWild v2 &
    CC BY-NC-ND 4.0 &
    No commercial use; no derivatives permitted.
    Wei et al., \textit{ACM MM}, 2024.
    DOI: 10.1145/3664647.3680599. &
    \smallurl{https://tqwei05.github.io/PlantWild/} \\
    \rowcolor{roweven}
    LeafBench &
    CC BY 4.0 &
    Released on Hugging Face as \texttt{enalis/LeafNet}; $\sim$70\% of data is public (training split only).
    Quoc et al., arXiv:2602.13662, 2026. &
    \smallurl{https://huggingface.co/datasets/enalis/LeafNet} \\
    \bottomrule
  \end{tabular}
\end{table*}

\endgroup


\section{Entropy Field Implementation}
\label{sec:app-entropy}

\paragraph{Extraction.} \texttt{logprobs=True} passed to vLLM client.
Per-token entropy: $h_i=-\log p_i$ (single-sample estimate, exact under
greedy decoding). $H_t^{(\mathrm{dis})}$ extracted from the token span
of PathogenAgent's \texttt{disease\_name} JSON field.

\paragraph{Hyperparameter sensitivity.}
$\delta_1\in\{0.10,0.15,0.20,0.25\}$ and $\delta_2$ tuned on 500-image
held-out set; full sensitivity table reported here after experiments.

\paragraph{Cross-model.} Entropy values comparable across Qwen3-VL-8B,
Phi-4-multimodal, Qwen2.5-VL-7B because all use same temperature and
the entropy field is read from the generation distribution, not
logit scaling.

\section{Budget Sensitivity}
\label{sec:app-budget}

\begin{table}[h]
  \centering\small
  \begin{tabular}{lccccc}
    \toprule
    $\beta$ & T3 F1 & T2 F1 & ECE & Mean $L$ & TPCP \\
    \midrule
    0 (no BT)   & --- & --- & --- & --- & --- \\
1 (default) & --- & --- & --- & --- & --- \\
2 (two BT)  & --- & --- & --- & --- & --- \\
\bottomrule
  \end{tabular}
  \caption{Backtrack budget sensitivity on PlantDoc (1k subset). Expected:
    $\beta{=}1$ optimal; $\beta{=}2$ marginal gain at $\sim$15\% token overhead.}
  \label{tab:budget}
\end{table}

\section{Path Clustering}
\label{sec:app-clustering}

Paths encoded as strings over \{Mo,Sy,Pa,Se,Di\}. Pairwise Levenshtein
distances feed hierarchical agglomerative clustering (Ward linkage);
cluster count maximizes silhouette score. Per-cluster disease-category
distribution and quality statistics characterize which image types trigger
which routing strategies.

\section{Hedge Lexicon}
\label{sec:app-hedge}

Seed terms: \textit{possibly, unclear, hard to tell, uncertain,
approximately, may, might, could, appears to be, cannot rule out}.
Co-occurring n-grams extracted within 5-token window, scored by PMI.
Top-200 PMI candidates manually reviewed. Validated against BioScope
\citep{vincze2008bioscope} and CoNLL 2010 Shared Task.

\end{document}